\chapter{Solution Methods}
\label{ch5}

The SA-VRPTW remains NP-hard after the
$\varepsilon$-constraint correction is imposed. This chapter
documents three solution approaches that all consume the same
\texttt{Instance} dataclass, all evaluate $F_1$ through
\texttt{savrptw.eval.objective.breakdown()}, and all satisfy or fail
the same \texttt{savrptw.eval.feasibility.validate()} contract. The
deep-reinforcement-learning solver discussed in earlier proposal
drafts has been removed because the current branch does not contain a
defensible, measured DRL baseline for the corrected formulation.

\section{Mixed-Integer Linear Programming}
\label{ch5:milp}

A PuLP/HiGHS-or-CBC model encodes the formulation in Chapter~3
directly. The implementation detail worth flagging is that two
arrival-variable groups are kept separate---$a_i^k$ for customer
arrivals and $a_{\text{return}}^k$ for the home-depot return---to
prevent the dispatch ($a_d^k = 0$) and return ($a_d^k > 0$)
constraints from contradicting each other for routes that visit at
least one customer. The convex SLA penalty is linearised by SOS2
outer-envelope segments at the breakpoints
$\tau \in \{0, 2, 5, 10, 15, 20, 30\}$~min listed in
Section~\ref{ch3:pwl}.

The default backend is HiGHS, with CBC as a free fallback when HiGHS
is unavailable. The solver is given a 120~s wall-clock budget per
instance in the default benchmark grid; longer budgets are needed
only for $N \geq 50$, where the branch-and-cut search is genuinely
difficult. Crucially, the big-$M$ in the arrival-linking
constraints~\eqref{eq:link} is set tight to $T_{\max}$ plus the
longest super-arc rather than to a hard-coded loose value, and CBC
presolve is disabled by default. Both choices were necessary to
eliminate spurious infeasibility observed during early calibration.

The MILP serves two roles: a quality benchmark on smaller instances
($N \leq 50$) and a correctness anchor for the heuristics. Where the
MILP returns a proven optimum, the heuristics' incumbents are
expected to lie within a small known gap.

\section{Genetic Algorithm}
\label{ch5:ga}

The genetic algorithm uses the giant-tour-with-Bellman--Ford-split
decoder of \citet{prins2004} embedded inside the standard genetic
operator framework \citep{holland1975}. Each chromosome is a
permutation over the customer set $C$. Decoding partitions the giant
tour into depot-respecting routes by solving a shortest-path problem
on an auxiliary acyclic graph whose arcs correspond to feasible
single-route cuts.

Operators are OX1 ordered crossover plus a randomised mutation choice
over swap, insertion, and segment reversal (the last subsumes the
2-opt move). Tournament selection uses $k = 5$, elitism retains the
top four individuals each generation, the default population size is
100, and the search runs for 200 generations.

The fitness evaluation uses $F_1$ plus a finite penalty mass for
infeasible routes. The penalty weights are tuned so that capacity,
duration, the per-route survival budget, and the per-route
residential cap each carry a stiff additive sentinel in the
Bellman--Ford split: a candidate split that violates any of those
caps is selected only if no feasible split exists for the same depot
sequence. This guarantees that, when the instance is feasible, the
GA returns a feasible incumbent.

\begin{pseudocode}{Giant-tour GA with Bellman--Ford split for SA-VRPTW.}
\label{alg:ga}
\textbf{Input:} \> Instance $\mathcal{I}$; population size $P$;
generations $G$; \\
\>\> crossover prob.\ $p_x$; mutation prob.\ $p_m$; \\
\>\> tournament size $k$; elite size $e$. \\
\textbf{Output:}\> Best feasible solution $s^\star$. \\[2pt]
\textbf{1.}\> Initialise population $\mathcal{P}_0$ with $P$ random
customer permutations. \\
\textbf{2.}\> Decode each chromosome via Bellman--Ford split; evaluate $F_1$ + penalty mass. \\
\textbf{3.}\> \textbf{for} $g = 1$ \textbf{to} $G$ \textbf{do} \\
\textbf{4.}\>\> Carry top-$e$ elites into next generation. \\
\textbf{5.}\>\> \textbf{while} offspring count $< P - e$ \textbf{do} \\
\textbf{6.}\>\>\> $p_1, p_2 \gets$ tournament selection of size $k$. \\
\textbf{7.}\>\>\> Child $c \gets$ OX1 crossover of $(p_1, p_2)$ with prob.\ $p_x$. \\
\textbf{8.}\>\>\> Apply random mutation (swap / insert / reverse) with prob.\ $p_m$. \\
\textbf{9.}\>\>\> Decode $c$ via Bellman--Ford split; evaluate fitness. \\
\textbf{10.}\>\> \textbf{end while} \\
\textbf{11.}\>\> Sort by fitness; truncate to $P$. \\
\textbf{12.}\> \textbf{end for} \\
\textbf{13.}\> Decode incumbent best chromosome; \\
\>\> validate strictly via \texttt{eval.feasibility.validate}. \\
\textbf{14.}\> \textbf{return} $s^\star$.
\end{pseudocode}

\paragraph{Bellman--Ford split.}
Given a chromosome $\sigma = (\sigma_1, \dots, \sigma_n)$, the
decoder computes, for every prefix index $j$, the cheapest cost
$F_1(0..j)$ to cover customers $\sigma_1..\sigma_j$ as a sequence of
routes ending at the home depot, by minimising over split points $i$
of $F_1(0..i) + \mathrm{routeCost}(\sigma_{i+1..j})$. Capacity bounds
the inner loop, so the decoder runs in $O(n^2)$ time. The route-cost
function shares the arrival-time recurrence with the MILP, so the GA
and MILP cannot disagree on a route's $F_1$.

\section{Adaptive Large Neighbourhood Search}
\label{ch5:alns}

The ALNS implementation follows \citet{ropke2006}. Destroy operators
include random removal, worst removal (largest insertion-cost
contribution), Shaw removal (relatedness over travel-time
proximity), and a custom risk-cluster removal that targets customers
in the highest-$R_{ij}$ neighbourhoods. Repair operators include
cheapest-feasible insertion, regret-2, and regret-3.

Operator weights are updated in segments of 100 iterations using the
Ropke--Pisinger reaction factor; rewards are
$\sigma = (33, 9, 1, 0)$ for new-best, improving, accepted-worse, and
rejected moves. Acceptance is simulated annealing with a geometric
cooling rate of $0.99975$, starting from a temperature equal to $5\%$
of the initial $F_1$.

\begin{pseudocode}{Adaptive Large Neighbourhood Search for SA-VRPTW.}
\label{alg:alns}
\textbf{Input:}\> Instance $\mathcal{I}$; iterations $T$; segment $L$; \\
\>\> reaction factor $r$; reward weights $\sigma$; \\
\>\> cooling rate $\eta$; init.\ temperature ratio $\theta$. \\
\textbf{Output:}\> Best feasible solution $s^\star$. \\[2pt]
\textbf{1.}\> $s_0 \gets \mathrm{GreedyConstruct}(\mathcal{I})$. \\
\textbf{2.}\> $s \gets s_0$;\ $s^\star \gets s_0$;\ $F^\star \gets F_1(s_0)$. \\
\textbf{3.}\> $T_{\mathrm{SA}} \gets \theta \cdot F^\star$. \\
\textbf{4.}\> $w^d_o, w^r_o \gets 1$ for every operator $o$. \\
\textbf{5.}\> \textbf{for} $t = 1$ \textbf{to} $T$ \textbf{do} \\
\textbf{6.}\>\> $q \gets \lceil n \cdot \mathcal{U}(0.15, 0.45) \rceil$. \\
\textbf{7.}\>\> $D \gets$ destroy operator sampled $\propto w^d$. \\
\textbf{8.}\>\> $R \gets$ repair operator sampled $\propto w^r$. \\
\textbf{9.}\>\> $s' \gets R(D(s, q))$;\ $\Delta \gets F_1(s') - F_1(s)$. \\
\textbf{10.}\>\> \textbf{if} $F_1(s') < F^\star$ \textbf{then} \\
\>\>\> $s^\star \gets s'$;\ $F^\star \gets F_1(s')$;\ $\rho \gets \sigma_1$. \\
\textbf{11.}\>\> \textbf{else if} $\Delta < 0$ \textbf{then} $\rho \gets \sigma_2$. \\
\textbf{12.}\>\> \textbf{else if} $\mathcal{U}(0,1) < \exp(-\Delta/T_{\mathrm{SA}})$
\textbf{then} $\rho \gets \sigma_3$. \\
\textbf{13.}\>\> \textbf{else} $\rho \gets \sigma_4$. \\
\textbf{14.}\>\> \textbf{if} $\rho > \sigma_4$ \textbf{then} $s \gets s'$. \\
\textbf{15.}\>\> $\pi^d_D \mathrel{+}= \rho$;\ $\pi^r_R \mathrel{+}= \rho$. \\
\textbf{16.}\>\> \textbf{if} $t \bmod L = 0$ \textbf{then} \\
\>\>\> $w^*_o \gets (1{-}r)\, w^*_o + r\, \pi^*_o / n^*_o$;\ reset $\pi^*$. \\
\textbf{17.}\>\> $T_{\mathrm{SA}} \gets \eta \cdot T_{\mathrm{SA}}$. \\
\textbf{18.}\> \textbf{return} $s^\star$.
\end{pseudocode}

\paragraph{Initial construction.}
The greedy constructor processes customers in random order per home
depot, appending each to the current route while feasibility holds.
When a singleton route is itself infeasible (e.g.\ the depot-to-customer
super-arc alone exceeds the residential cap), the customer is kept as
a singleton anyway and the SA loop is given a chance to relocate it.
The strict \texttt{validate()} call at the end remains the final
feasibility gate, so genuinely infeasible instances surface as a
clear infeasibility error rather than as a silent failure.

\section{Crash-Survival Monte Carlo}

For evaluation only (not optimisation), each route is sampled
$10\,000$ times: each edge traversal is a Bernoulli trial with
success probability $1 - r_e$, and the route survives iff every edge
traversal succeeds. The reported summary contains mean route survival,
fleet survival, and Wilson 95\% confidence intervals
\citep{wilson1927}. The analytic per-route survival
$\exp(-R_{\text{route}})$ is reported alongside the Monte Carlo
estimate as a sanity check.

\section{Behavioural Compliance Sensitivity}
\label{ch5:behav}

Rider compliance is modelled as a $\mathrm{Beta}(\alpha, \beta)$
prior. The central scenario is $(\alpha, \beta) = (5, 2)$. The
sensitivity sweep covers $\alpha \in \{3, 5, 8\}$ and
$\beta \in \{1, 2, 3\}$, yielding nine operating points. For each
point, the evaluator Monte-Carlos a 50-rider population over 5{,}000
trips and reports the inter-quartile-range band on per-rider
survival. Edge-level traversal rates are inflated by a factor
$1 + k_c (1 - c)$, where $c$ is the sampled compliance level and
$k_c = 0.4$ is matched to the Indian gig-rider field data of
\citet{das2020}. The behavioural evaluator is a diagnostic layer
only; it does not modify the optimisation objective or the solver
comparison itself.

\section{Comparative Summary}

Table~\ref{tab:solver_comparison} summarises the role each solver
plays in the corrected architecture.

\begin{table}[ht]
\centering
\caption{Comparative roles of the three solution approaches.}
\label{tab:solver_comparison}
\small
\begin{tabular}{llll}
\toprule
\textbf{Method} & \textbf{Role} & \textbf{Strength} & \textbf{Limitation} \\
\midrule
MILP & Exact baseline ($N \leq 50$) & Formulation fidelity & Limited scalability \\
GA   & Fast heuristic     & Strong giant-tour sequencing & Sensitive to penalty calibration \\
ALNS & Large-instance heuristic & Powerful repair under tight budgets & Operator-tuning sensitivity \\
\bottomrule
\end{tabular}
\end{table}

The three solvers are kept inside a single $F_1$/feasibility contract.
This is the methodological move that makes apples-to-apples
comparison possible across exact and metaheuristic results.
